\documentclass[10pt,a4paper]{article}
\usepackage[utf8]{inputenc}
%\usepackage[german]{babel}
%\usepackage[cyrillic]{babel}
\usepackage[T1]{fontenc}
\usepackage{amsmath}
\usepackage{amsfonts}
\usepackage{amssymb}
\usepackage{amsthm}
%\usepackage{stix-type1}%\turnediota
\theoremstyle{definition}
\newtheorem{definition}{Definition}
\newtheorem{lemma}{Lemma}
\newtheorem{theorem}{Theorem}
\usepackage{etoolbox}

%\newcommand*{\@Dashv}[2]{%
%  \reflectbox{$\m@th#1#2$}%
%}
\newcommand*{\separation}[3]{% 1)binder, 2)set, 3)condition
  \{#1\in #2: #3\}
}

% https://www.freecodecamp.org/news/simple-regex-tricks-for-beginners-3acb3fa257cb/
%  \\mathcal\{T\} |---> \topology
%  \\mathcal\{([^\}]*)\}   |--->  $1
\newcommand{\set}{\mathrm{\ set}}
\newcommand{\univocal}{\mathrm{\ univocal}}
\newcommand{\singleton}{\mathrm{\ singleton}}
\newcommand{\complex}{\mathbb{C}}
\newcommand{\real}{\mathbb{R}}
\newcommand{\power}{\mathcal{P}}
\newcommand{\borel}{\mathcal{B}}
\newcommand{\interior}{\mbox{int}}
\newcommand{\closure}{\mbox{cl}}
\newcommand{\nghbset}{\mathcal{N}}
\newcommand{\semialgebra}{\mathcal{S}}
\newcommand{\topology}{\mathcal{T}}
\newcommand{\union}{\bigcup}
\newcommand{\functional}{\mbox{\ functional}}
\newcommand{\injective}{\mbox{ }}
\newcommand{\total}{\mbox{\ total on }}
\newcommand{\surjective}{\mbox{\ surjective on}}
\newcommand{\countable}{\mbox{\ countable}}
\newcommand{\converges}{\longrightarrow}
\newcommand{\AND}{\ \&\ }
\newcommand{\IMPLIES}{\ \Rightarrow\ }
\newcommand{\Proof}{\noindent \textbf{Proof.}\\}
\newcommand{\Done}{\noindent \textbf{Done.}\\}
\newcommand{\UB}{\mathrm{UB}}
\newcommand{\LB}{\mathrm{LB}}
\newcommand{\Sup}{\mathrm{Sup}}
\newcommand{\Inf}{\mathrm{Inf}}
% local commands
%\preto\section{\undef\blipo} %etoolbox
\preto\subsection{\undef\blipo} %etoolbox
% https://tex.stackexchange.com/questions/105684/limit-macro-newcommand-scope-to-section
\begin{document}
%qwertyuio
\noindent

The aim of this document is to provide perfect proofs of functional analysis and topology. It is needed because writing it by hand leads to severe amount of corrections, while implementing it with an interactive proof assistant is tedious and platform-dependent.\\
%\noindent
% elements are small letters, sets are uppercase latin letters, sets of sets are mathcal latin letters
% the code should be comprehensive: use meaning of symbols
\noindent
\section{Goals}
Interesting targets:
\begin{itemize}
\item sequential criterion
\item completion of a metric space
\end{itemize}
\newpage
\section{Logic}
$\mathcal{A}$ Alphabet



\subsection{Syntax}
Basic notion of the system is "class".
\subsubsection{Terms}
Terms ::= \\

\subsubsection{Formulas}
\begin{definition}
Fm ::=\\
\end{definition}

\begin{definition}
$$M(Z) \Leftrightarrow \exists C: Z\in C$$
\end{definition}

G\"odel notation:\\

\subsection{Axioms}
interpretation is denoted by I here.

Usually assignment is an univocal correspondence between variables and sets.
In our case assignment is an univocal correspondence between variables and conglomerate-singletons of classes.\\

There are several possible ways of interpreting "in" relation: 
Assume X and Y are some terms(sic! not just variables).
1) Weak
$X\in Y$ is valid iff for any assignment any class-element of I(X) is in any class-element of I(Y).\\

Hypothesis:
$$A(\varepsilon x: A)\Leftrightarrow \exists x: A(x)$$




Axiom schema



Following class existence metatheorem for NBG, we can introduce an axiom schema.\\
To begin with let's introduce recursive function on formulas: it substitutes class variables with set variables.


The resulting formula contain only set variables.



If $\varphi$ contains only set variables, then 
$\{x | \varphi\}$ is defined.
$$(\varphi\in\mathrm{setify}[\mathrm{Fm}])\stackrel{meta}{\Longrightarrow} ( \vdash (\{x | \varphi\}\downarrow)) $$

\subsection{Semantics}
"super-large" model.\\
$K$ - conglomerate of all classes.\\

Assignment:\\
$$A:\mbox{Var}\to \power(K)$$

$$I^{\mathrm{Fm}}_A(S\downarrow) \Leftrightarrow (I^{\mathrm{Term}}_A(S)\neq\emptyset)$$

Definiton of truth:
$$\mbox{Bool} := \{\mbox{False}, \mbox{True}\}$$
$$M,A\vDash \varphi \stackrel{meta}{\Longleftrightarrow} I_A(\varphi)=\mbox{True}$$

$$(M\vDash \varphi) \stackrel{meta}{\Longleftrightarrow} (\forall A\in \mbox{Assignment}: M,A\vDash \varphi)$$

%I_A(\varphi)=\mbox{True}
\subsection{Extension by definitions}


\subsection{Example of a proof (FARMER'S CALCULUS)}
Partially based on WILLIAM M. FARMER \& JOSHUA D. GUTTMAN's article\\
``A Set Theory with Support for Partial Functions''.\\
For simplicity of understanding by the reader, the notation from the article is used here.\\
\begin{definition} Value of a function.
$$f(x) \simeq \big(\iota y: \langle x,y\rangle\in f\big)$$
\end{definition}

\begin{theorem} If $f$ is univocal, then
$$\big(f(x)\simeq y\big)\Leftrightarrow \big(\langle x,y\rangle\in f\big)$$
\end{theorem}
\Proof
($\Leftarrow$) (Obvious part) Here we use axiom PFOL$10'$ which is equaivalent to PFOL10:
$$
(\iota x:\varphi)\downarrow
\quad\overset{\mbox{PFOL10}'}{\Longleftrightarrow}\quad
(\exists!x:\varphi)
$$
So the proof is the following:
$$
\big(\langle x,y\rangle\in f\big)
\quad\overset{\exists\mbox{-intro}}{\Longrightarrow}\quad
\big(\exists w: \langle x,w\rangle\in f\big)
$$
$$
f\mbox{ univocal}
\quad\overset{\mbox{def}}{\Longleftrightarrow}\quad
\big(! w: \langle x,w\rangle\in f\big)
$$
$$
\big(\exists w: \langle x,w\rangle\in f\big)
\quad\&\quad
\big(! w: \langle x,w\rangle\in f\big)
\quad\overset{\exists!\mbox{-intro}}{\Longrightarrow}\quad
\big(\exists! w: \langle x,w\rangle\in f\big)
$$
$$
\big(\exists! w: \langle x,w\rangle\in f\big)
\quad\overset{\mbox{PFOL10}'}{\Longrightarrow}\quad
\big(\iota w: \langle x,w\rangle\in f\big)\downarrow
$$
$$
\big(\iota w: \langle x,w\rangle\in f\big)\downarrow 
\quad\overset{\mbox{PFOL11}}{\Longrightarrow}\quad
\langle x,(\iota w: \langle x,w\rangle\in f)\rangle\in f
$$
$$
\langle x,(\iota w: \langle x,w\rangle\in f)\rangle\in f
\quad\&\quad
\big(\langle x,y\rangle\in f\big)
\quad\overset{f \mbox{ univocal}}{\Longrightarrow}\quad
(\iota w: \langle x,w\rangle\in f) \simeq y
$$
$$
f(x) \simeq y
$$
%\big(! w: \langle x,w\rangle\in f\big)
%$$
%\big(\iota w: \langle x,w\rangle\in f\big)\downarrow

%$$
\noindent($\Rightarrow$)
$$(y\downarrow)\quad\mbox{(by PFOL8)}$$
$$
(y\downarrow)\quad\&\quad((\iota w: \langle x,w\rangle\in f)\simeq y) \quad\overset{\mbox{PFOL7}}{\Longrightarrow}\quad
(\iota w: \langle x,w\rangle\in f)\downarrow $$
$$
\big(\iota w: \langle x,w\rangle\in f\big)\downarrow 
\quad\overset{\mbox{PFOL11}}{\Longrightarrow}\quad
\langle x,(\iota w: \langle x,w\rangle\in f)\rangle\in f
$$
$$
\langle x,(\iota w: \langle x,w\rangle\in f)\rangle\in f
\quad\&\quad
\big((\iota w: \langle x,w\rangle\in f)\simeq y\big)
\quad\overset{\mbox{PFOL7}}{\Longrightarrow}\quad
\big(\langle x, y\rangle\in f\big)
$$
\\
\Done
\pagebreak
\subsection{Example of a proof (MORSE-KELLEY SET THEORY, CLASSICAL SEMANTICS)}

\begin{definition} Simultaneous substitution.
$$\varphi[Z/W]$$
\end{definition}

\begin{definition} Set.
$$(X\set)\ \Leftrightarrow\ (\exists Y:X\in Y)$$
\end{definition}

\begin{definition} Class.
$$Y\in\{x|\varphi[x/W]\} \Leftrightarrow (Y\set)\ \&\ (\varphi[Y/W])$$
%$$Y\in\{x|\varphi\} \Leftrightarrow (Y\set)\ \&\ (\varphi[Y/x])$$
%(or maybe
%$$Y\in\{x|\varphi[x/Y]\} \Leftrightarrow (Y\set)\ \&\ (\varphi)$$
%we need to find one expression for both cases.
%)
\end{definition}


\begin{definition} Universal class.
$$V=\{x|\top\}$$
\end{definition}


\begin{theorem}
$$(X\set)\ \Leftrightarrow\ (X\in V)$$
\end{theorem}
\Proof
\noindent($\Leftarrow$) Assume $(X\in V)$.
$$(X\in V)\ \Longrightarrow\ (\exists W: X\in W).$$
$$(\exists W: X\in W)\ \Longrightarrow\ (X\set).$$
\noindent($\Rightarrow$) Assume $(X\set)$.
%$$(X\set)\ \Longrightarrow\ (\exists W: X\in W))$$
$$(X\set)\ \&\ \top[X/W]\ \Longrightarrow\ (X\in \{x|\top[x/W]\})$$
$$(X\in \{x|\top[x/W]\}=\{x|\top\}=V)\ \Longrightarrow\ x\in V$$
%$$(a\in\{a\})\ \&\ (\{a\}=\{w\})\ \Longrightarrow\ (a\in\{w\})$$
\Done
\begin{theorem}
$$A=B\ \Longrightarrow\ (t[A/W]=t[B/W])$$
\end{theorem}
\Proof
By induction on term construction...\\
\Done

\begin{theorem}
$$A=B\ \Longrightarrow\ (\varphi[A/W]\ \Longleftrightarrow\ \varphi[B/W])$$
\end{theorem}
\Proof
By induction on formula construction...\\
\Done

\begin{definition} Epsilon operator.
$$(\varepsilon y: \varphi) = c\big(\{y|\varphi\}\big)$$
\end{definition}

\begin{definition} Iota operator.
$$(\iota y: \varphi) = \bigcap \{y|\varphi \land (!Y:\varphi)\}$$
\end{definition}

\begin{definition} Lambda operator.
$$(\lambda x: t) = \{p|\exists x: p = \langle x,t \rangle\}$$
\end{definition}


%\Leftrightarrow
\begin{definition} Value of a function.
$$F(x) = \big(\iota y: \langle x,y\rangle\in F\big)$$
\end{definition}


\begin{definition} Singleton.
$$(A\singleton) \Longleftrightarrow (\exists!a:a\in A)$$
\end{definition}

\begin{theorem}
$$(w\in A)\land (!a:a\in A) \longrightarrow \{w\}=A$$
\end{theorem}
\Proof
\noindent($\subseteq$) Assume $(q\in \{w\})$.
$$(q\in \{w\})\ \Longrightarrow\ (q=w).$$
$$(w\in A)\ \&\ (q=w)\ \Longrightarrow\ (q\in A).$$
\noindent($\supseteq$) Assume $(a\in A)$.
$$(!a:a\in A)\ \&\ (a\in A)\ \&\ (w\in A)\ \Longrightarrow\ (a=w).$$
$$(a=w)\ \Longrightarrow\ (\{a\}=\{w\})$$
$$(a\in\{a\})\ \&\ (\{a\}=\{w\})\ \Longrightarrow\ (a\in\{w\})$$
\Done

\begin{theorem}
$$\big(\bigcap\{w\}=w\big)$$
\end{theorem}
\Proof
\Done
\begin{theorem}
$$(A \singleton) \longrightarrow \big(\big(\bigcap A\big)\in A\big)$$
\end{theorem}
\Proof
$$(A\singleton)\ \Longrightarrow\ (\exists!w:w\in A)$$
$$(w\in A)\ \&\ (!a:a\in A)\ \Longrightarrow\ \{w\}=A\quad\mbox{(by the previous theorem)}$$
$$(\{w\}=A)\ \Longrightarrow\ \big(\bigcap\{w\}=\bigcap A\big).$$
$$\big(\bigcap\{w\}=w\big)\ \&\ \big(\bigcap\{w\}=\bigcap A\big)\ \Longrightarrow\ \big(\bigcap A=w\big)$$
$$\big(\bigcap A=w\big)\ \&\ (w\in A)\ \Longrightarrow\ \big(\bigcap A\in A\big)$$
\Done



\begin{theorem} If $F$ is univocal, then
$$\big(F(x)=y\big)\ \Longleftrightarrow\ \big(\langle x,y\rangle\in F\big)$$
\end{theorem}
\Proof
($\Leftarrow$) Assume $(\langle x,y\rangle\in F)$.
$$.$$
($\Rightarrow$) Assume $(F(x)=y)$.
$$(F(x)=y)\ \Longrightarrow\ \big(\iota y: \langle x,y\rangle\in F\big)=y.$$
$$\big(\iota y: \langle x,y\rangle\in F\big)=y\ \Longrightarrow\ \big(\bigcap\{y: \langle x,y\rangle\in F\}\big)=y.$$
$$(F\univocal)
\ \Longrightarrow\ 
(\forall x:\ !y : \langle x,y\rangle \in F)
$$
$$
(\forall x:\ !y : \langle x,y\rangle \in F)
\ \Longrightarrow\ 
(!y : \langle x,y\rangle \in F)
$$
$$
(y\set)
\ \&\ 
\big(\bigcap\{y: \langle x,y\rangle\in F\}\big)=y
\ \Longrightarrow\ 
(\big(\bigcap\{y: \langle x,y\rangle\in F\}\big)\set)
$$

$$
\big(\{y: \langle x,y\rangle\in F\}\singleton\big)
\ \Longrightarrow\ 
\big(\bigcap\{y: \langle x,y\rangle\in F\}\big)
\in
\big(\bigcap\{y: \langle x,y\rangle\in F\}\big)\quad(\mbox{by ...}).$$
\Done

\section{Singleton-Semantics}
\subsection{•}
$$((\varepsilon x:x\in Q)\rightsquigarrow t) \longleftrightarrow (t\in Q)$$

$$((\varepsilon y:\psi)\rightsquigarrow t) \longleftrightarrow \psi[t/y]$$
\section{Set \& class theory}
Notation:\\

Starting here and for the rest of the document, we will use usual symbol of equality "=" for quasi-equality (before it was "$\simeq$") and "$\equiv$" for equality (before it was "=").

\begin{definition} Identity.
$$I = \{(x,x)|\}$$
\end{definition}

\begin{definition} "F on X".
$$\mbox{"F on X"} = F\restriction X $$
\end{definition}


\begin{definition}
$$y\in f[A] \Leftrightarrow \exists x\in A: \langle x,y\rangle\in f$$
\end{definition}
\begin{definition}
$$y\in f\langle A\rangle \Leftrightarrow \forall x\in A: \langle x,y\rangle\in f$$
\end{definition}
\begin{theorem}
$$X\subseteq \power Y \Leftrightarrow \union X \subseteq Y$$
\end{theorem}
\begin{theorem}
$$g\functional \AND M \countable \IMPLIES g[M] \countable$$
\end{theorem}
\begin{theorem}
$$f^\leftarrow \functional$$
\end{theorem}

\begin{theorem}
$$f\injective \IMPLIES f\left[\union M\right] = \union f^{\rightarrow}\left[M\right]$$
\end{theorem}

\begin{theorem}
$$f\functional \IMPLIES f^{-1}\left[\union M\right] =\union f^{\leftarrow}\left[M\right]$$
\end{theorem}
\section{Linear algebra}
\begin{definition}Matrix over a set S.
$$\mbox{"A is a matrix over a set S"} \Leftrightarrow \exists m, n \in \mathbb{N}: \{1,\dots,m\}\times \{1,\dots,n\} \to S $$
\end{definition}

\begin{definition}
$$\mbox{combine}(f_1,f_2) := \{(((m,k),(m,v)),(m,(k,v)))| (m,k)\in f_1 \land (m,v)\in f_2\}$$
\end{definition}
(A combination is a product of morphisms in the category sets and functions.)

\section{Real analysis}
\begin{definition}
$$\UB_R(S) = R[S]$$
\end{definition}
\begin{definition}
$$\LB_R(S) = R^{-1}[S]$$
\end{definition}

Context and notation:\\
Here by $\leq$ we denote the standard relation "less or equal to" on real numbers.
$\leq$ for inequalities of numbers\\
$\leqslant$ for inequalities of sets of numbers\\
$R := \leq$\\
$\leq = \geq^{-1}$\\
$\geq = \leq^{-1}$\\

\noindent

\begin{definition}
$$A \leqslant B \Longleftrightarrow \forall a\in A \forall b\in B: a\leq b$$
\end{definition}
\begin{definition}
$$\UB(S) = \leq [S]$$
\end{definition}
\begin{definition}
$$\LB(S) = \geq [S]$$
\end{definition}

\begin{theorem}
$A\leqslant B \Longleftrightarrow B\subseteq \UB(A)$\\
\end{theorem}
\begin{theorem}
$A\leqslant B \Longleftrightarrow A\subseteq \LB(B)$\\
\end{theorem}

\begin{definition}
$$\Sup(S) = \UB(S)\cap \LB(\UB(S))$$
\end{definition}
\begin{definition}
$$\Inf(S) = \LB(S)\cap \UB(\LB(S))$$
\end{definition}


\begin{definition}
$$\mbox{sup}(S) = (\iota x: x\in \mathrm{Sup}(S))$$
\end{definition}
\begin{definition}
$$\mbox{inf}(S) = (\iota x: x\in \mathrm{Inf}(S))$$
\end{definition}


\begin{theorem}
$$A \leqslant B \Longrightarrow A\leqslant \Sup(A)\leqslant \Inf(B)\leqslant B$$
\end{theorem}


\begin{theorem}
$$a:\mathbb{N}\rightarrow \real, \{a_n\}_n \mbox{bdd above in }\real, (a_n)_n\mbox{ non-decreasing} \Rightarrow a \converges \sup \{a_n\}_n$$
\end{theorem}
\section{Topology}
\begin{definition}
$$\mbox{Segm} = \{[-\infty;\beta):\beta\in\real\}\cup \{(\alpha;+\infty]:\alpha\in\real\}\cup \{(\alpha;\beta):\alpha\in\real,\beta\in\real\}$$
\end{definition}
\begin{definition}
$$\topology_{[-\infty;+\infty]} = \mbox{Segm}_\sigma$$
\end{definition}
\begin{definition}
Set of neighbourhoods
$$\nghbset_\topology(x) = \{S \subseteq \union \topology : \exists  U\in \topology: x\in U \subseteq S \}$$
\end{definition}
\begin{definition}
Interior of a set V in a topology $\topology$.
$$\interior_\topology(V) = \union
\separation{W}{\topology}{W\subseteq V}$$
\end{definition}
\begin{theorem}
$$(\interior = \nghbset^{-1})\AND(\nghbset=\interior^{-1})$$
\end{theorem}
\begin{theorem}
$$V\in \nghbset_\topology\langle A\rangle \Leftrightarrow A \subseteq \interior(V) $$
\end{theorem}
\begin{theorem}
$$\closure \dashv \mbox{id} \dashv \interior$$
\end{theorem}

\newcommand*{\blipo}{\separation{W}{\topology}{W\subseteq V}}
% \{W\in\topology|W\subseteq V\}
%\noindent Lemma:
\begin{theorem}
$$V\in\nghbset_\topology(x) \Leftrightarrow x\in \interior_\topology(V)$$
\end{theorem}
\Proof
$(\Rightarrow)\\
V\in \nghbset_\topology(x) \IMPLIES \exists U\in \topology: x\in U \subseteq V$\hfill(nghb-elim, union-elim).
\\
$U\in\topology\AND U\subseteq V \IMPLIES U \in \blipo $\hfill(specification-intro).
\\
$x\in U \AND U \in \blipo \IMPLIES x\in \union \blipo$\hfill(union-intro).
\\
$x\in \union \blipo \IMPLIES x\in \interior_\topology(V)$\hfill(interior-intro).
\\
$(\Leftarrow)$
\\
$x\in \interior_\topology(V) \IMPLIES \exists W\in  \{U\in\topology :U\subseteq V\}:\ x\in W$\hfill(interior-elim, union-elim).
\\
$W\in  \{U\in\topology:U\subseteq V\} \IMPLIES W\in\topology\AND W\subseteq V$\hfill(specification-elim).
\\
$x\in W\AND W\in\topology\AND W\subseteq V \IMPLIES \exists U\in \topology: x\in U \subseteq V$\hfill(and-intro, exists-intro).
\\
$(\exists U\in \topology: x\in U \subseteq V) \IMPLIES V\in \nghbset_\topology(x)$\hfill(exists-elim, nghb-intro).
\\
\Done
\section{Measure theory}
\section{Functional analysis}
% $\blipo$ should give error
\section{Probability theory}
\section{Category theory}
\end{document}\\
